\section{怎样进行记叙文的缩写、扩写或改写}

《中学语文教学大纲》明确指出：作文训练的方式是多种多样的，缩写、扩写、改写等等，同命题作文一样，也是训练和考查学生的写作能力的重要方式。为什么这样说呢？因为不论缩写、扩写或改写，都属于“按照材料作文”的类型，写作时，先得“读”材料，然后才是“写”作文，这样既免除了习作者挖空心思编造材料之苦，又培养了习作者的读、写能力。所以教育界老前辈叶圣陶称赞这种“按照材料作文”的写作方式“打破了命题作文的老传统，是思想上的大突破，大解放。”

缩写、扩写或改写，尽管议论文、说明文、应用文也可以采取这些方式进行写作训练，但是都不如记叙文运用得广泛而多样，因此我们对记叙文的缩写、扩写或改写，更要加以注意并经常练习。

怎样进行记叙文的缩写、扩写或改写呢？

\textbf{1. 缩写的要点是删除枝叶，保留主干，缩得精要。}
所谓缩写，就是把内容丰富的长文章压缩成为言简意赅的短文。这种作文训练方式的特点是：一定要保留原文的主要思想内容，而略述或删削那些属于旁枝末节的部分，达到既突出了原文的基本情节，又显得短而不空的写作要求，从而培养习作者从纷繁复杂的长文中理清脉络、抓住要点、组织语句、扼要表达的能力。具体说来，在缩写记叙文时应该注意：

（1）首先必须认真阅读所供的材料，深入理解所供材料的思想内容，把握所供材料的中心思想，然后围绕这个中心思想，找出各段的内容要点，并进一步划出表示内容要点的关键语句，进行合理的剪裁，有机地组织材料，从而列出缩写的段落提纲或要点，再考虑下笔行文。这是缩写成败的关键，因为找错了要点或者遗漏了要点，就不能反映原文的精神实质和动人风貌，甚至将轻重倒置，歪曲原貌。

（2）缩写的重点只在“缩”字上，因此必须忠实于原文，不要随意改变原文的中心和要点，更不要随意改变原文的体裁、人称和写作方法。

（3）在找准原文的主要情节、主要人物之后，要在缩写过程中把它们作为“主干”突出出来，而把原文中的次要情节、次要人物以及某些说明成分、交代成分等作为“枝叶”进行略写或删去，并把原文中细致的描写、具体的叙述进行概括。但是，这种压缩和概括必须做到“压”而不“瘪”，“概”而不“空”，不能使缩写后的文章过于抽象，只剩几条筋，成了原文的提纲或片断摘录；当然更不能随意发挥自己的看法，写成原文的读后感或分析评论。

（4）缩写后的记叙文，在结构上要注意保持原文的连贯性和完整性。为了语气贯通，前后衔接，缩写时可以适当添加过渡性的语句。

（5）在语言上可以尽量保留原文中关键的和十分精彩的语句，而对那些在原文中属于铺陈、描述、联想、举例、发挥、人物对话等部分则应略写，或者用概括转述的方式捎带一笔，甚至在不妨碍文意连贯、有损要点的前提下可以全部删去。缩写后的语言风格应该努力与原文保持一致。

例如缩写法国作家都德的著名小说《最后一课》（初中语文课本第二册）。通过阅读和分析，我们知道这是一篇表现法国人民爱国思想的小说。作者以普法战争为背景，通过对“最后一课”的描写，刻画了小弗朗士、韩麦尔等性格鲜明的人物形象，揭示出国土沦丧对法国人民的巨大震撼。当我们明确文章的中心思想之后，就可以围绕中心思想来确定文章的主要情节、主要人物。很明显，这篇小说的全部情节可以分成课前、课上和下课三个部分。课前部分是交代普法战争法国失败的背景，并由普鲁士占领军贴出布告为下文故事的展开打下伏笔，这不是故事的主要情节，只要略写出“我”上学迟到而无心看布告就行了。课上部分是“最后一课”的主要情节，而韩麦尔先生所说的“这是我最后一次给你们上课了”则是点题之笔，并推动了故事的发展，那是应该照写的关键句。上课的内容既有法语课，又有习字课，和历史课，而重点则在“我的最后一堂法语课！”法语课的重点又在韩麦尔先生没有责备小弗朗士背不出书，只是讲了必须把法国语言记在心上的那一段有名的话。那么我们就突出缩写上法语课的情景，其它则略写或删去。下课部分是故事的尾声，要点在于韩麦尔先生使出全身的力量写出的“法兰西万岁”，我们突出这一点就够了。至于人物，“我”是不能少的，因为他是贯串全文的人物，然而比起韩麦尔来，还不能作为主要的描写对象，至于铁匠华希特、郝叟老头、从前的镇长、邮递员等次要人物只能略写一笔，也不必一一提出名字了。在这样找出主要情节、主要人物和关键语句之后，就可以动手编写缩写提纲，也就有可能把这篇两千多字的小说用六、七百字“精粹”地表述出来了。你如果有兴趣，不妨缩写一下试试看。

\textbf{2. 扩写的要点是紧扣关键，添枝加叶，扩得合理。}

所谓扩写，就是将所供的材料扩展成内容丰富多采的文章。这种作文训练方式的特点是：必须在不离原意的基础上，扩展和扩充必要的合理的内容。“扩展”是说依据原文展开描写，进行润饰，“扩充”是说依据原文充实材料，丰富内容。“扩展”又不“越轨”，扩写后的文章与原文在一般道上“跑车”，“扩充”又不累赘，扩写后的文章，不是原文的“虚胖”。进行这样的作文练习，能开拓学生的思路，培养学生的想象能力、表达能力和逻辑推理能力。具体说来，扩写记叙文时应该注意：

（1）必须首先认真钻研原文，把握其中心思想，围绕其中心思想理清原文的脉络，确定其主要情节、主要人物，找出文中的关键的话，作为扩写时应该着力发挥的依据。

（2）扩写的重点只在“扩”字上。我们要在忠实于原文的主要内容的基础上，充分发挥想象力，按照合理的逻辑和生活的规律展开原文，使原文内容更充实，细节更生动。但是，扩展的地方，只是对原文关键处的精雕细刻，而不要一味堆砌词藻，也不要把短句都变成长句，更不要随意渲染与中心思想关系不大的语句，扩充的地方，只是对原文关键处的合理补充，而不要脱离原文节外生枝，凭空虚构情节，使补充的情节游离于原文中心之外，成为文章的累赘。

（3）原文的主要情节、主要人物，扩写时都不能随意丢掉、增加或更换；原文表述的形式也不能随意改变。对原文情节的扩写，必须合乎原文内容和生活规律；对人物形象的扩写，必须合乎原文的人物性格和思想感情。

（4）扩写时，一定要注意文章的结构，做到思路清晰，脉络贯通，结构完整，层次分明。一般不要改变原文的记叙顺序，也不要由于扩展了一些合情合理的描绘或添加了一些必要的细节，而使文章冗长、拖沓，结构松散。

（5）扩写过程中，还要注意多种表现手法和修辞方法的运用，使文章更具生动性，但切忌凭空添上许多不着边际的议论、抒情或说明的话语。

例如秦牧的著名散文《土地》（见于高中语文课本第二册）中关于晋公子重耳出亡途中接受一个农民给予的土块的描述，就是对《左传·僖公二十三年》中的一段有关记载的扩写。《左传》的有关记载只是说：“晋公子重耳过卫，出于五鹿，乞食于野人，野人与之块。公子怒，欲鞭之。子犯曰：‘天赐也。’稽首，受而载之。”但在秦牧的笔下，这一段极为简炼的文字却成了一个有头有尾、丰富多采的故事：

一队亡命贵族，在黄土高原上仆仆奔驰。他们虽然仗剑驾车，然而看得出来，他们疲倦极了，饥饿极了。他们用搜索的眼光望着田野，然而骄阳在上，田垄间麦苗稀疏，哪里有什么可吃的东西！一个农民正在田里除草。那流亡队伍中一个王子模样的人，走下车子来，尽量客气地向农民请求：“求你给我们弄点吃的东西吧！你总得帮忙才好。我们已经好几天没有吃的了。”衣不蔽体、家里正在愁吃愁穿的农民望了望这群不知稼穑艰难的人们一眼，一句话也没说，从田地里捧起一大块泥土，送到王子模样的人面前，压抑着悲忿说：“这个给你吧！”王子模样的人显然被激怒了，他转身到车子上去拿马鞭，怒气冲冲地想逞一下威风，鞭打那个胆敢冒犯他的尊严的农民。但是一个上了年纪的、大臣模样的人上前去劝阻了：“这是土地，上天赐给我们的，可不正是我们的好征兆么！”于是，一幕怪剧出现了，那王子模样的人突然跪下，叩头谢过苍天，然后郑重地捧起土块，放在车上，一行人又策马前进了。辘轳大车过处，卷起了漫天尘土······

请看，秦牧同志扩写后的文字，完全保留了原文的主要情节、主要人物，也没有改变原文的思想内容和记叙形式，但是，他从原文的内容和形式出发，进行了合乎生活规律和合乎原文人物性格的想象：首先他根据原文中“过卫，出于五鹿”的记载，合理想象了公子重耳等在黄土高原上奔驰逃亡的情景，添枝加叶地对骄阳在上、麦苗细疏作了细节描写，实在是在“意料之外，情理之中”，为下文农民没有给食物而给泥块作了铺垫。其次他根据原文中“乞食于野人，野人与之块”的记载，合理想象了公子重耳怎样尽量客气地“乞”食，而农民怎样饥寒交迫、愤而给予土块的情景，把原文客观的叙述扩展为具体生动的描写，使得比较抽象的叙述变成了活灵活现的图景；再次，他根据原文中“公子怒，欲鞭之。子犯曰：‘天赐也。’”的记载，合理想象了公子重耳这个逃亡王子由他的阶级本性所决定了的蛮横行为和子犯这个多谋大臣一心想保全公子重耳重新夺取晋国王位的机智话语，扩展得完全符合生活逻辑和人物性格；最后他又根据原文中“稽首，受而载之”的记载，合理想象了“一幕怪剧”的情景，不惜笔墨地扩充了细节，对这一关键处精雕细刻，描绘得使人如临其境，如见其人。在结构上，这个故事扩写得又是多么完整、谨严，显得条理清楚，衔接紧密，前呼后应，首尾圆合。我们应从这样精彩的扩写文字里，细心体会扩写的一般要求和具体做法。

\textbf{3. 改写的要点是认准要求，重新表述，改得恰当。}

所谓改写，就是根据题目要求而改变原文的表现形式，对原文进行再创作的一种作文训练方式。这种作文训练方式的特点是：必须根据题目的具体要求，在表现形式（包括体裁、结构、人称、语言、记叙手法、表现角度等）的某些方面作相应的改变，使改写后的文章与原文“貌”离而“神”合，即表现形式变了而思想内容基本一致。

改写与缩写、扩写有一致的地方，就是都要符合原文的思想内容。但是，根据改写题目的具体要求，改写后的文章与原文相比，篇幅可长可短（作为作文考试的题目，大多要求改写后的文章比原文短），这跟缩写后的文章一定比原文短，扩写后的文章一定比原文长是不同的，特别是改写的重点就在“改”字上，它在表现形式方面总是要求对原文有所改动，而缩写和扩写在表现形式方面一般要跟原文相同。

改写的常见方式有：

（1）改变原文的体裁：如把诗歌改写为散文，把小说改写为剧本，把剧本改写为故事等。

（2）改变原文的结构：如把倒叙或插叙改写为顺叙，或者反过来把顺叙改为倒叙或插叙等。

（3）改变原文的人称：如把原文第一人称改为第三人称，或者把第三人称改为第一人称等。

（4）改变原文的表现手法：如把人物对话改为直接叙述，把侧面描写改为正面描写，把明线改为暗线，把主观想象改为客观记叙，把单纯的记叙改为夹叙夹议等。

（5）改变原文的写作角度：如把几个人物的故事改写为一个人物的故事，或者把次要人物改写成主要人物，把记人为主改为记事为主等。

改写时应该注意：

（1）首先要细心阅读原文，认真体会改写题目的要求，根据作文题目的要求要毫不含糊地作相应的改变；有时题目要求比较含蓄，更要仔细琢磨，努力弄清其内在的含义，认准有关表现方法的基本特点，然后才能下笔编写提纲。

（2）要在仔细阅读原文的基础上，确定改写的中心内容，分清主次，重要的部分要详写，次要的部分要略写，有积极意义的地方要着力写，有消极意义的地方则要避免写，以使内容更集中，主线更分明，中心更突出。

（3）改写一般可以根据题目的要求，对原文的材料进行合理的取舍，重新安排文章的结构，重新考虑记叙的手法，使文章具有改写的特色。

（4）改写应根据题目的要求，对体裁或人称、角度等作相应的改变，而一般不应在内容上作很大的改动或扩展，不能编造原文没有的新情节（丰富细节是允许的）。

（5）改写可以说是根据原文重写，进行再创作，因此一定要避免象缩写那样大量抄写原句、原段，力求用自己的语言表现原文的内容。

例如一九七九年全国高考语文试卷中的作文题，就是要求把何为同志所写的《第二次考试》改写成一篇《陈伊玲的故事》。题目明确要求：“按原文内容写一篇以陈伊玲为中心的记叙文，不要另外编造情节，不要写成《第二次考试》的缩写”，还要求：“要有明确的中心思想，注意材料的剪裁和组织”，“层次清楚，结构完整”，“语言通顺，标点正确，不写错别字”，“字数以六、七百字为好，最多不得超过八百字”等。我们可以根据这些要求，明确题目规定了改写的文体必须是“故事”性质的记叙文，而不能写成诗歌、读后感之类；规定了改写的中心人物是陈伊玲，而不能象原文那样将陈伊玲与苏林教授同作故事中心人物，从而改变了写作的角度；规定了改写的写法只能根据原文进行再创作，既不能另外编造情节，也不能象缩写那样可以整句、整段地照抄原文。弄清了题目的要求之后，就得仔细阅读原文，从而确定改写的中心内容：那就是歌颂陈伊玲不仅有优美的音色，更有高尚的品格，一定能成为一个优秀的歌唱家。根据这个中心内容，就能比较容易地选取原文中能够表现这一中心内容的材料：初试——音色优美，感情深沉，成绩出众；复试——声音发涩，毫无光彩，令人惊诧；原因——复试前夕，她整夜投入了扑灭火灾、安置灾民的工作，过分劳累；结果——主考人员弄清原委，破格录取了她。这四个方面的材料，又以连夜抢险救灾最为重要，因为它不仅是整个故事波澜起伏的关键，也是歌颂陈伊玲具有人民歌唱家高尚品格的闪光点，因此这个方面一定要突出写。第二次考试是原文的重点场面，改写时可以详点写；第一次考试是陪衬，可以略点些，结局的部分只要点明陈伊玲由于具备人民歌唱家的条件而被破格录取即可。至于原文中其他的材料大半可以舍去，甚至为了突出中心人物陈伊玲，苏林教授的大名也可以略去，只提“主考人员”或“老教授”也就可以了。为了体现改写的特色，结构上还应该重新组织。例如：先介绍陈伊玲的爱好音乐，更热爱工作；再写她初试出众；接着写她连夜救灾；然后写复试失利，最后被破格录取。这是一种打破原文结构、全部采取顺叙的写法。又如：先写陈伊玲不顾一切地连夜救灾，再写她参加第二次考试；接写主考人员对她第一次考试的回忆；最后是破格录取，这也打破了原文结构而采取了倒叙、补叙的写法。只要能做到层次清楚，结构完整，不管采用哪种结构方法都是可以的。另外，改写所用的语言除了通顺之外，还得表现出自己口语的特色，如果只是抄写原文中的这几方面的文字，那岂不成为缩写了吗？
\newpage